\documentclass[11pt]{article}

\usepackage[margin=1.05in]{geometry}
\usepackage{amsmath,amssymb,amsthm}
\usepackage{booktabs}
\usepackage{microtype}
\usepackage{xcolor}
\usepackage[colorlinks=true,linkcolor=blue!60!black,citecolor=blue!60!black,urlcolor=blue!60!black]{hyperref}

\newcommand{\PT}{\mathcal{PT}}
\newcommand{\ee}{\mathrm{e}}
\newcommand{\ii}{\mathrm{i}}
\newcommand{\norm}[1]{\left\lVert #1 \right\rVert}


\title{The multiplicity horn survives renormalization:\\
exact quantum evolution of a ghost--normal field theory,\\
the folded channel sum, and an accumulated-truncation floor}
\author{Jonathan Torkelson\\[2pt]
\small Embeddetech, Inc., Tulsa, Oklahoma, USA \quad
\texttt{jtorkelson@embeddetech.com}}
\date{July 13, 2026}

\begin{document}
\maketitle

\begin{abstract}
\noindent Two earlier papers of this program price the
ghost-sector trilemma on a premise: that the quantum vacuum-decay
rate of a coupled ghost sector scales with the number of open
resonant channels---the \emph{multiplicity horn}---so that the
divergent field-theoretic rate is a channel-counting statement,
not an artifact of perturbation theory. Against this stands the
possibility that structural effects beyond perturbation
theory---the position we test, framed as the
Gross--Strumia--Teresi--Zirilli (GSTZ) metastability
suggestion---could erode the channel sum. We put the premise
to a falsification test in a genuinely quantum, multimode
1+1-dimensional two-field theory (normal $\varphi$ plus ghost
$\chi$, contact coupling $\lambda\varphi^2\chi^2$, equal masses:
maximal resonance, with the four-particle
$\varphi\varphi\chi\chi$ shell exactly degenerate with the
vacuum) by exact sparse-propagator evolution and a static-$E_0$ Feshbach/Schur
folding of the full Hamiltonian onto the resonant shell. The
pre-registered fork read: (a) any deficit of the exact rate
against the bare channel sum is absorbed by standard
renormalization (dressing) of the states entering the golden rule, or (b) a structural
suppression survives folding. The data return (a), to the ${\sim}1\%$
level (with a documented downward drift of the static-$E_0$ fold):
across couplings where the \emph{naive} channel sum misses by
up to $31\%$, and across mode numbers up to $M=16$ where the
naive ratio slides to $0.774$, the folded (Schur) prediction
tracks the exact evolution at $0.988$--$1.000$. By
elimination---diagonal self-energies and vacuum-column vertex
corrections each reproduce the naive ratio to the four digits resolved---the
deficit is attributed to the shell--shell block: first-order splitting
and coherent mixing of the degenerate shell, textbook degenerate-level
renormalization, no new
mechanism. A dimension ladder calibrated to the exact Hamiltonian
to $5.7\times10^{-16}$ reproduces the program's reduced-measure
rungs in quantum dress: the exactly resonant $d=1$ shell
saturates with cutoff, the $d=1$ golden-rule sum grows log-class, and the
$d=2$ sum grows as a power (fitted exponent $2.26$ against the
predicted $\approx2$; a short-lever cutoff sweep, a mechanism check rather than a
scaling measurement). A matrix-product-state cross-check
certifies the position-space theory against exact diagonalization at $L=4$
(to $6\times10^{-6}$), runs in production at $L=16$, and then meets an honest
limit that we characterize rather than hide: the
Trotter/truncation floor is non-monotonic in the step
($0.337/0.057/0.231$ at $dt=0.04/0.02/0.01$, $D=40$), and
raising $D$ from $40$ to $64$ moves the floor from $0.057$
to $0.062$---saturation, per the pre-registered decision rule. The
volume-scaling extension is therefore left open at the MPS level;
the verdict itself rests on the exact methods. Conformal-gravity
pricing that leans on the multiplicity premise survives this
test; positions that require a nonperturbative erosion of the
channel sum do not receive support at any scale we can reach.
\end{abstract}

\newpage
\begin{quote}
\small\noindent\textbf{Note on provenance.} This manuscript is an
experiment in AI-conducted research as much as a physics paper. The
project began when the author asked the large language model Claude
Fable~5 (Anthropic): \emph{``Is it not possible that the antilinear
PT/time-reversal symmetry element could maybe predict indivisible
non-Markovian processes?''} Everything that followed was performed by
the model: the research program, the theorems and proofs, the numerical
experiments, the literature search, and this text. The model was
repeatedly prompted to attack its own results adversarially and to
search for prior work. The author, a professional engineer rather than
a physicist, directed and reviewed the sessions but makes no claim to
comprehension of the findings, and is circulating the manuscript for
independent review by professional physicists. Accordingly, nothing
here should be taken on trust: every numerical claim is reproducible
from the public scripts (see Code availability); the originating
conversation, project briefing, and full research notebook are
preserved in the program's hub repository,
\url{https://github.com/embeddetech/ghosts-program}.
\end{quote}

%=====================================================================
\section{Introduction}
\label{sec:intro}

This is the seventh paper of a program
that reads quantum dynamics in Barandes' currency of indivisible
stochastic processes on
configurations~\cite{Barandes,indivisible,instantaneous,conditioned,choosing,unfolded}. Two of its predecessors lean on one
premise. The vacuum-decay analysis of the ghost
sector~\cite{instantaneous} factorizes the golden-rule rate into a
boost-orbit volume times a reduced integral whose scaling is set
by channel counting; the conformal-gravity pricing
of~\cite{choosing} then prices the Lorentz-violating-cutoff horn
of the trilemma on that reduced measure. Both inherit the
field-theoretic position that the divergent decay rate of a
coupled ghost vacuum is a statement about the \emph{multiplicity}
of open resonant channels---what those papers call the
multiplicity horn.

That premise has a published antagonist. Gross, Strumia, Teresi,
and Zirilli, asking whether negative kinetic energy is even
metastable, concede the perturbative rate is ``naively infinite''
while suggesting the true nonperturbative rate is unknown and
possibly benign~\cite{GSTZ2021}; a recent line of work constructs
interacting ghost systems whose structure arrests the cascade
entirely~\cite{Deffayet2026,EwasiukProfumo2026}. If some
structural mechanism---beyond the textbook renormalization of the
states entering the golden rule---eroded the channel sum in a
genuinely quantum evolution, the premise under two papers of this
program would fail, and with it their pricing of the trilemma.

This paper is the falsification test. We build the smallest
genuinely quantum, multimode theory in which the question is
sharp: a normal field $\varphi$ and a ghost field $\chi$ in
$1{+}1$ dimensions, momentum modes on a circle, contact
interaction $\lambda\int\varphi^2\chi^2$, equal masses. The ghost
sign makes the four-particle $\varphi\varphi\chi\chi$ states with
balanced momenta \emph{exactly degenerate} with the vacuum: the
resonant shell of the golden rule is not an approximation here
but an exact zero-energy subspace, and its population count is
the channel multiplicity. The pre-registered fork
(recorded in the program notebook before the decisive runs):
\begin{itemize}
\item[(a)] the exact decay of the vacuum tracks the channel sum
after standard renormalization of the states entering it---the
multiplicity premise survives; or
\item[(b)] a deficit survives the folding---incipient structural
suppression, and the premise fails.
\end{itemize}
The data return (a) at every scale we can reach
(Sec.~\ref{sec:fold}). The deficit that perturbation theory does
show is attributed, by elimination, to first-order splitting and
coherent mixing of the degenerate shell
(Sec.~\ref{sec:mechanism})---renormalization, not suppression. A
dimension ladder reproduces the reduced-measure rungs
of~\cite{instantaneous} in quantum dress (Sec.~\ref{sec:ladder}). A
matrix-product-state cross-check certifies the position-space
theory against ED at $L=4$ (production at $L=16$) and then hits a truncation floor that we
measure and report as a methods limit rather than extrapolate
through---an accumulated-truncation floor (Sec.~\ref{sec:mps}).

Honest scope, stated at the outset: everything here is
$1{+}1$-dimensional (with a $d=2$ spot check), at most $M=16$
($32$ modes per field), occupation caps to $10$, and a static-energy Schur fold.
What the test shows is that finite-channel quantum dynamics
contains no mechanism eroding the channel sum beyond textbook
renormalization \emph{above the ${\sim}1\%$ level reached here}; the
infinite-channel, fully resummed limit---the
true boost-orbit volume---remains analytic territory, priced
in~\cite{instantaneous} and not re-derived here.

%=====================================================================
\section{The model and the channel census}
\label{sec:model}

On $2M$ momentum modes per field ($k_n=2\pi n/\ell$,
$n\in\{\pm1,\dots,\pm M\}$, zero mode dropped, on a circle of length
$\ell$; the small-$M$ stage-3 checks, the Sec.~\ref{sec:model} channel
ratios, the $d=2$ spot check, and the ladder calibration use $\ell=8$,
while the campaign, rate-window, born-vacuum, and dimension-ladder
scans use $\ell=16$),
with equal masses ($m=1$) and free Hamiltonian
$H_0=H_\varphi-H_\chi$ (the ghost enters with the Ostrogradski
sign), the contact interaction
$\lambda\int\varphi^2\chi^2$ couples the Fock vacuum to
four-particle states $\varphi_{k_1}\varphi_{k_2}\chi_{k_3}
\chi_{k_4}$ with $\sum k_i=0$. At equal masses the subset with
$\{k_1,k_2\}=\{-k_3,-k_4\}$ pairings has free energy exactly
zero: the resonant shell. The initial state is the
negative-norm-free Fock vacuum of the free theory, whose
excitations carry negative energy---the arena of the
GSTZ framing~\cite{GSTZ2021}; we note that the integral-of-motion
construction of~\cite{Deffayet2026} is built on $(x^2-y^2)$-type interactions and
does not extend to a pure $\varphi^2\chi^2$ coupling (there the biquadratic
enters only as an explicitly nonintegrable perturbation), so their
mechanism does not protect this coupling.

The channel census (\texttt{scripts/t5\_channel\_ladder.py}) counts
resonant channels $10/21/36/55$ at $M=2/3/4/5$ (the coupled
resonant pairings, already including the same-mode pairs; the full
zero-energy subspace $P$ quoted in Table~\ref{tab:verdict} is larger,
since it also collects the zero-energy two-particle
$\varphi_k\chi_{-k}$ pairs---and, at larger particle-number caps,
higher-number degenerate states---not only these four-particle
pairings);
channel couplings
are $g=\lambda/(\ell\sqrt{w_1w_2w_3w_4})$ (the $w_i$ the four
participating mode frequencies) with a $1/\sqrt2$ per
identical pair, and the analytic channel table matches the exact
Hamiltonian's vacuum column to $5.7\times10^{-16}$
(\texttt{scripts/t5\_dimension\_ladder.py})---the calibration that anchors
the ladder to the same theory. The archived check is at $d=1$,
$\ell=8$, $M=4/6$, cap~4; the $d=2$ rung applies the same
box-normalized vertex without an archived exact-Hamiltonian
comparison of its own.

Exact sparse-propagator evolution of $\langle N_\varphi\rangle(t)$ against
the first-order channel sum gives ratios $0.999$--$0.995$ out to
$t=20$ at $\lambda=0.02$: at weak coupling the naive
multiplicity premise is already nearly exact. The question is
what happens when it is not.

%=====================================================================
\section{The decisive test: folding the full theory onto the
shell}
\label{sec:fold}

At larger coupling the naive channel sum overshoots the exact
evolution. The decisive question is whether the deficit is
\emph{renormalization} (the states entering the golden rule are
dressed) or \emph{suppression} (something erodes the
multiplicity itself). The instrument that separates them is a
(static-$E_0$) Feshbach/Schur fold: with $P$ the projector onto the vacuum
plus the full degenerate shell and $Q$ the rest,
\[
H_{\rm eff}\;=\;PHP\;-\;PHQ\,(QHQ-E_0)^{-1}QHP,
\qquad E_0=0,
\]
evolved exactly within $P$. The fold keeps every channel---the
multiplicity is untouched---while dressing the shell and the
vacuum by the full $Q$-space to second order in the fold and to
all orders in the $Q$-space dynamics. If the deficit is
renormalization, the folded prediction tracks the exact
evolution; if multiplicity erodes by a mechanism beyond the reach of the
static second-order dressing, it cannot.

\begin{table}[t]
\centering\small
\begin{tabular}{@{}llll@{}}
\toprule
configuration & full/naive & full/Schur & \\
\midrule
$M=2/3/4$, cap 4, $\lambda=0.02$, $t=40$ & --- &
$1.0000\,/\,0.9997\,/\,0.9993$ & \\
$M=2/3/4$, cap 6 & --- & $0.9998\,/\,0.9984\,/\,0.9963$ & \\
$M=4$, cap 4, $\lambda=0.04$ & $0.9434$ & $0.9972$ & \\
$M=10$, cap 4, $\lambda=0.02$ & $0.8959$ & $0.9951$ & \\
$M=12$, cap 4 & --- & $0.9941$ & \\
$M=16$, cap 4 & $0.7737$ & $0.9924$ &
{\footnotesize($|P|=561$, $|Q|=15{,}416$)}\\
$M=6$, cap 6 & --- & $0.9884$ &
{\footnotesize($|Q|=22{,}056$)}\\
\bottomrule
\end{tabular}
\caption{Exact evolution versus the naive channel sum and the
Schur-folded channel sum
(\texttt{scripts/t5\_stage3.py}, \texttt{scripts/t5\_campaign.py}). Where the
naive sum misses by $23\%$ ($M=16$), the folded sum tracks the
exact rate to $0.8\%$. The $M\le4$ rows are the $\ell=8$ stage-3
checks and the $M\ge6$ rows the $\ell=16$ campaign ($m=1$ throughout);
the naive ratio is box-dependent at the $\sim1$--$2$ percentage-point
level, but the Schur-folded band is robust to $\ell$, so the verdict
does not depend on the box.}
\label{tab:verdict}
\end{table}

Table~\ref{tab:verdict} is the verdict. Across the campaign the
folded ratio sits in the band $0.988$--$1.000$ while the naive
ratio slides from $0.995$ to $0.774$. The $\lambda$-scaling of
the naive deficit at $M=3$, cap 6 ($0.0057/0.0227/0.0883/0.3144$ across
$\lambda=0.005\to0.04$, close to $\times4$ per doubling)
identifies its leading piece as $\lambda^2$---the $(\Delta E\,t)^2$
signature of an $O(\lambda)$ splitting of the exactly degenerate
shell---with the sub-$\times4$ growth at the top of the range
marking a negative (saturating) higher-order ($\lambda^4$) correction; the folded column stays at
$0.0001$--$0.0069$ throughout, including where the naive sum
misses by $31\%$. Occupation-cap convergence (caps $4/6/8/10$ at
$M=2$, dimensions to $2568$) holds the folded ratio flat at
$0.9998$--$1.0000$ while the naive ratio settles
non-monotonically near ${\sim}0.956$: the deficit is not a truncation artifact, and
neither is its absorption by the fold. A $d=2$ spot check (the
$48$-channel shell $\{-1,0,1\}^2\setminus0$, cap 4) gives
full/Schur $=1.0000$: the same behavior holds in the one $d=2$
shell we can reach.

One honest sentence the table earns: the folded ratio drifts
mildly downward as the fold deepens ($1.0000$ at small $M$ to
$0.9884$ at the largest $|Q|$ we can invert), which is
consistent with the expected degradation of the \emph{static}-$E_0$ Schur
approximation rather than a resurgent deficit (fork (b) is thereby
excluded only above the ${\sim}1\%$ level at the deepest reachable
fold)---but we state it rather
than absorb it, and the full fold at $|Q|\approx81{,}000$
(cap-6 $M=8$) exceeded the memory available to this program; the
record there is the elimination chain of
Sec.~\ref{sec:mechanism}, not a folded ratio.

%=====================================================================
\section{What the deficit is: attribution by elimination}
\label{sec:mechanism}

At $M=8$, cap 6, the naive ratio is $0.8620$. Two partial
dressings were computed exactly
(\texttt{scripts/t5\_partial\_folds.py}) and both \emph{fail} to move it:
folding only the diagonal self-energies reproduces the naive
ratio exactly, and folding the vacuum-column vertex correction
reproduces it exactly again ($0.8620$ both times). What remains---%
and what the full fold captures at reachable $M$---is the
\emph{shell--shell block}: the direct couplings among degenerate shell
members, their $Q$-induced second-order corrections, and the vertex
corrections to non-vacuum columns. Both partial dressings retain the
naive first-order \emph{structure} (a channel-by-channel sum of
$|g_f|^{2}$ terms with no shell--shell rotation), which is why neither
can move the ratio. Because the free energies vanish on the shell, the
shell--shell block splits it at first order in $\lambda$; entering the
transition probability as $(\Delta E\,t)^{2}$, that splitting reproduces
the $\lambda^{2}$ growth of the deficit measured in
Sec.~\ref{sec:fold}. No wave-function $Z$-factor is involved: the static
fold is Hermitized and evolved unitarily inside $P$
(\texttt{scripts/t5\_stage3.py}), so it carries $Z\equiv1$ by
construction, and in the exact evolution the \emph{total} norm leaked
out of $P$ at this configuration is $6.3\times10^{-3}$ against a
$13.8\%$ naive deficit---over an order of magnitude too small to be the
leading piece. The deficit is thus attributed, by elimination
(each partial dressing individually null; a shell-shell-\emph{only} fold was not
separately computed, and the full fold at this
configuration is memory-limited, so the attribution leans on the
reachable-$M$ folds of Table~\ref{tab:verdict}), to degenerate-shell
splitting---dressed kinematics, not eroded multiplicity. (The notebook records a failed expectation here:
the author predicted the vertex correction would carry the
deficit at the ${\sim}0.98$ level; it carried none of it. The
prediction and its failure are both logged, per the program's
pre-registration discipline.)

Two further signatures close the dynamical story
(\texttt{scripts/t5\_stage4.py}). The two-component window for the
vacuum-survival rate $\Gamma(t)$ (the log-derivative of the survival
probability), $\Gamma(t)=a+bt$ at cap~4, returns $b=4.28\times10^{-5}$ against
$2S_{\rm exact}=4.53\times10^{-5}$ (${\sim}6\%$; $S_{\rm exact}=\sum_{\rm shell}|g|^{2}$, each exactly-resonant $E\!\approx\!0$ channel's squared coupling---the matrix-element modulus squared computed in the script), with the constant
term $a=2.03\times10^{-4}$ sitting just above the golden-rule plateau
($1.56$--$1.84\times10^{-4}$). And the born-vacuum signature---in a deliberately detuned twin
($m_{\chi}=1.35$, where the pair-channel detunings die as $1/k$ in the
UV so the energy--time window $|\Delta E|<2\pi/t$ can close; the twin is
detuned in the \emph{spectrum} only---the vertex normalization retains
the $m=1$ mode functions, a smooth $O(1)$ reweighting immaterial to the
window kinematics read off here): the
transfer rate is \emph{sustained then collapsing} while the number of
channels inside that window falls ($471\to7$ across $t=5\to120$)---the
vacuum decays as a multichannel resonance while channels are open and
stops when they close. Whether the rate itself grows logarithmically with
time in the born-vacuum regime is not resolvable at $M=8$ and we
do not claim it; the sustained-then-collapse profile is the
signature shown.

%=====================================================================
\section{The dimension ladder, in quantum dress}
\label{sec:ladder}

The reduced-measure analysis of~\cite{instantaneous} assigns
vacuum-decay integrals to rungs by spacetime dimension and vertex
class. With the channel table calibrated to the exact Hamiltonian
($5.7\times10^{-16}$, Sec.~\ref{sec:model}), the same rungs can
be read quantum mechanically. Three pre-registered predictions,
all confirmed (\texttt{scripts/t5\_dimension\_ladder.py}; the
golden-rule sums use a Gaussian smoothing width $\eta=0.1$, and the
quoted exponents shift by ${\sim}0.01$--$0.15$ across
$\eta=0.05$--$0.2$ without changing any verdict):
the exactly resonant $d=1$ shell \emph{saturates} with cutoff
(P1: $+1.5\%$ over the last doubling at $\Lambda=100$); the
$d=1$ golden-rule sum grows \emph{log-class}, the channel-table
image of Garriga--Vilenkin's $D=2$ logarithm~\cite{GV2013} (P2:
power-fit
exponent $0.156$ with still-decaying increments---consistent with
a logarithm approaching its asymptote, and stated as log-class
evidence rather than a converged asymptote); and the $d=2$ sum
grows as a \emph{power} (P3: fitted exponent $2.263$ against the
predicted ${\approx}2$; the $d=2$ scan is a short-lever cutoff sweep
(cutoff $M=4\text{--}8$, three-point fit), so this
is a mechanism check, not a scaling measurement). Channel-table
extensivity holds within ${\sim}7.4\%$ under volume doubling. (The
notebook also records that the author initially mislabeled which
dimension carries the log rung; the docstring was corrected
against~\cite{instantaneous} before the first run.)

%=====================================================================
\section{The MPS cross-check, and the accumulated-truncation floor}
\label{sec:mps}

The exact methods above live in momentum space. As a fourth,
independent cross-check we evolve the position-space lattice
theory (two bosons per site, open boundaries,
$H=H_\varphi-H_\chi+\lambda\sum_i\varphi_i^2\chi_i^2$) by
second-order TEBD, preparing the free vacuum as the
imaginary-time ground state of the sign-flipped auxiliary
Hamiltonian $H_\varphi+H_\chi$ (the same state; the ghost sign
enters real-time evolution only), and reading the normal-sector
energy transfer $\Delta E_\varphi(t)$ per interior site.

\emph{Certification.} The gate run against dense exact
diagonalization at $L=4$ ($D=100$, $dt=0.02$, second-order
Trotter, initialized from the exact ground state via dense
diagonalization of the auxiliary Hamiltonian) agrees to
$2$--$6\times10^{-6}$ at $t=1\ldots4$---up to three orders below
signal (by $t=4$; ${\sim}50\times$ below at $t=1$). At $L=16$, $D=40$, $dt=0.02$ (the production
protocol): the $\lambda=0$ floor is $\max|\Delta E|=5.73\times
10^{-2}$ ($7.16\times10^{-3}$ per interior site), the
$\lambda=0.06$ normal-sign twin sits at $5.41\times10^{-2}$, and
the $\lambda=0.06$ ghost run shows $\Delta E_\varphi/{\rm site}
=8.93\times10^{-3}$ by $t=12$---a signal $24.8\%$ above the floor.
At $D=40$ only the $\lambda=0.12$ rung separates cleanly from
the floor (${\sim}1.9\times$, on the diagnostic $dt=0.01$ ladder
described below); the $\lambda=0.03$ rung is
floor-buried, so the coupling-scaling claim of this paper is
carried by the exact methods of Secs.~\ref{sec:fold}
and~\ref{sec:ladder}, not by an MPS $\lambda$-ladder.

\emph{The floor, measured.} The floor is not a precision
issue: at $10^{-2}$--$10^{-1}$ it sits fourteen orders above
machine epsilon, and it is non-monotonic in the Trotter step---%
$0.337/0.057/0.231$ at $dt=0.04/0.02/0.01$ ($D=40$,
$L=16$)---the classic signature of a Trotter/truncation tradeoff:
Trotter error falls as $T\,dt^2$ while truncation error, the
Schmidt weight discarded at each re-compression to $D$,
accumulates with the step count ${\sim}T/dt$. The sum is
minimized near $dt=0.02$, and refining the step \emph{backfires}
(a shared expectation of author and director, refuted by the
data and logged). The mechanism is \emph{not} a free-field entanglement
barrier. At $\lambda=0$ the free vacuum is an exact eigenstate of the
real-time generator, and the $L=4$ gate prepares it exactly (dense
diagonalization of the auxiliary Hamiltonian); there its middle-cut
entropy across the middle bond is stationary---$0.2740$ at every
$t\in[0,12]$---so there is no
free-field entanglement
growth to outrun the $S\sim\log D$ capacity. Nor is the floor static
representability alone: at $L=4$ the exact vacuum's discarded Schmidt
weight beyond rank $40$ at that bond is $2.2\times10^{-11}$, far below the observed
floor. What accumulates is Trotter-injected excitation, repeatedly
re-truncated: in the production protocol ($L=16$, $D=40$, $dt=0.02$) the
\emph{simulated} state's middle-cut entropy climbs from $0.249$ at $t=0$
to $0.89$ at $t=12$, whereas at $L=4$---the only size at which an exact
reference exists---the exact state does not move and the $D=40$
simulation tracks it to ${\sim}10^{-3}$ ($0.2740\to0.2752$). The
$L=16$ runs prepare the vacuum variationally, by annealed
imaginary-time TEBD at $D=40$, so their climb mixes Trotter-injected
excitation with imperfect eigenstate preparation; the $L=4$, $D=100$
control (flat to $5\times10^{-5}$) is what shows the climb is not
intrinsic. The pre-registered decision run
($D=64$, $dt=0.02$, $L=16$) returned a floor of
$6.22\times10^{-2}$ against $D=40$'s $5.73\times10^{-2}$: a
$60\%$ increase in bond dimension (a $13\%$ increase in
$\log D$ capacity) did not lower the floor (it rose slightly). Per the decision
rule fixed before the run, we conclude the MPS is
truncation-limited at this coupling and report it as a methods
limit: \emph{no volume-scaling claim is made from MPS data}. The
diagnostic $dt=0.01$ ladder (including an $L=32$ point whose
per-site transfer came out negative, i.e.\ non-extensive under
truncation) is exactly that---diagnostic; volume extensivity at
the field level is left open at the MPS level, supported here
only by the exact channel-table extensivity (${\sim}7.4\%$) and the
$L=4$ ED gate test.

For an exact integrator none of this arises---which is precisely
why the verdict of this paper rests on ED and the Schur fold,
with the MPS supplying an independent $L=4$ certification (production at $L=16$) and an
honestly characterized boundary of applicability.

%=====================================================================
\section{Relation to prior work}
\label{sec:prior}

Real-time false-vacuum decay on quantum lattices is an active
lane: spin-chain false-vacuum decay~\cite{Lagnese2021}, qumode
MPS simulations of scalar-field decay---the nearest
methodological template to Sec.~\ref{sec:mps}%
~\cite{AbelSpannowskyWilliams2025}, and Fermi-golden-rule onset
studies in discrete-mode systems~\cite{BurshteinGoldstein2024}.
On the ghost side, the state of the art in exact quantum
treatments is the exact operator-conservation-law moment bounds
of~\cite{EwasiukProfumo2026} and the integral-of-motion
construction of~\cite{Deffayet2026}; the metastability suggestion
tested here is that of~\cite{GSTZ2021}, and the channel-counting
frame descends from Cline--Jeon--Moore~\cite{CJM2004}; the low-dimensional log
structure enters via the reduced-measure ladder of~\cite{instantaneous}. To our knowledge no
published work evolves a multimode ghost--normal theory in real
time, and none tests vacuum decay against channel count; both
statements were re-checked against the 2025--2026 literature at
drafting. The result is
load-bearing in opposite directions for two camps: for the
structural-stability program~\cite{Deffayet2026,EwasiukProfumo2026} it marks where structure must enter (their
mechanisms are interaction-structural, absent for
$\varphi^2\chi^2$); for the metastability
suggestion~\cite{GSTZ2021} it removes the only
finite-channel mechanism we could construct.

%=====================================================================
\section{Discussion}
\label{sec:discussion}

What is established, at what strength. \emph{Exact}: the
four-particle shell's exact degeneracy with the vacuum by
construction, and the analytic channel-table calibration (matching the
exact vacuum column to $5.7\times10^{-16}$). \emph{Empirical
(static-$E_0$ fold)}: the folded channel sum tracks the exact sparse-propagator
evolution within the band $0.988$--$1.000$ across every configuration
where the fold is computable (Table~\ref{tab:verdict}), with the documented
downward drift of the static-$E_0$ approximation; the two null partial
folds of Sec.~\ref{sec:mechanism} (the shell-splitting attribution they
support is inferred, not exact).
\emph{Computed, pre-registered}: the dimension-ladder rungs (P1
saturation; P2 log-class evidence, asymptote not pinned; P3
power, short-lever cutoff sweep). \emph{Certified}: the MPS propagator against ED at $L=4$, $D=100$
(production at $L=16$, $D=40$), with a measured truncation floor beyond it. \emph{Not
claimed}: volume scaling from MPS; the born-vacuum log-growth of
the rate; the infinite-channel resummed limit, which remains the
analytic content of~\cite{instantaneous}.

For the program, the consequence is simple: the multiplicity
premise under the vacuum-decay pricing of~\cite{instantaneous} and
the conformal-gravity fork of~\cite{choosing} survives its own
falsification test. A structural suppression of the channel sum,
if one exists, must live where this test cannot reach---in the
infinite-channel limit, in higher dimensions at depth, or in
interaction structures engineered for
stability~\cite{Deffayet2026,EwasiukProfumo2026}---and not in the
generic finite-channel quantum dynamics of a coupled ghost.

\paragraph{Code availability.} All numerical claims are
reproducible from the public scripts in the program repository:
{\emergencystretch=2em\def\_{\char`\_\allowbreak}\catcode`\/=13 \def/{\char`\/\allowbreak}%
\texttt{scripts/t5\_channel\_ladder.py}, \texttt{scripts/t5\_stage2.py},
\texttt{scripts/t5\_stage3.py}, \texttt{scripts/t5\_stage4.py},
\texttt{scripts/t5\_dimension\_ladder.py}, \texttt{scripts/t5\_campaign.py},
\texttt{scripts/t5\_partial\_folds.py}
(Secs.~\ref{sec:model}--\ref{sec:ladder});
\texttt{scripts/t5\_mps.py}, \texttt{scripts/t5\_mps\_gate.py},
\texttt{scripts/t5\_mps\_prod.py}, \texttt{scripts/t5\_floor\_diag.py},
\texttt{scripts/t5\_run\_one.py}, \texttt{scripts/t5\_prod\_dt01.py},
\texttt{scripts/t5\_floor\_mechanism.py} (Sec.~\ref{sec:mps}; the last
emits the entropy, Schmidt-weight and norm-leak diagnostics, and prints
the entropy at \emph{every} bond so the quoted middle-bond values carry
their cut); a numerics-provenance
map \texttt{scripts/SS5\_NUMERICS.md}; result files
(\texttt{results\_val\_M*\_cap*.txt},
\texttt{results\_partial\_folds\_M8\_cap6.txt},
\texttt{mps\_gate\_result.txt}, \texttt{mps\_prod\_results.txt},
\texttt{dt01\_*.txt}, \texttt{diag\_*\_c64\_dt02.txt}) archived
alongside.\par}

\begin{thebibliography}{99}\small

\bibitem{indivisible}
J.~Torkelson,
\emph{Indivisible stochastic processes for non-Hermitian and
higher-derivative quantum systems: detailed balance, broken $\PT$
symmetry, and the explosion of ghosts},
repository \url{https://github.com/embeddetech/p01-indivisible-ghosts}
(2026).

\bibitem{instantaneous}
J.~Torkelson,
\emph{Ghost vacuum decay as a stochastic process: instantaneous
states, the boost-orbit measure, and the foliation that ghosts make
visible}, repository
\url{https://github.com/embeddetech/p03-instantaneous-ghosts} (2026).

\bibitem{conditioned}
J.~Torkelson,
\emph{The price of a post-selected signal: a sharp
survival--distinguishability trade-off, saturated by broken-$\PT$
dynamics}, repository
\url{https://github.com/embeddetech/p04-conditioned-ghosts} (2026).

\bibitem{choosing}
J.~Torkelson,
\emph{Conformal gravity must choose: exceptional-point kinematics,
vacuum decay, and the sample space of Weyl gravity}, repository
\url{https://github.com/embeddetech/p05-choosing-ghosts} (2026).

\bibitem{unfolded}
J.~Torkelson,
\emph{The conformal point is a phase boundary: unfolding
directions, the sign map of the cascade divisor, and a two-channel
photonic readout}, repository
\url{https://github.com/embeddetech/p06-unfolded-ghosts} (2026).

\bibitem{Barandes}
J.~A.~Barandes,
\emph{The stochastic-quantum correspondence},
Philosophy of Physics \textbf{3}(1), 8 (2025),
doi:10.31389/pop.186, arXiv:2302.10778;
\emph{The stochastic-quantum theorem}, arXiv:2309.03085.

\bibitem{GSTZ2021}
C.~Gross, A.~Strumia, D.~Teresi, and M.~Zirilli,
\emph{Is negative kinetic energy metastable?},
Phys.\ Rev.\ D \textbf{103}, 115025 (2021), arXiv:2007.05541.

\bibitem{Deffayet2026}
C.~Deffayet, A.~Fathe Jalali, A.~Held, S.~Mukohyama, and
A.~Vikman,
\emph{Unitary time evolution and vacuum for a quantum stable
ghost},
arXiv:2604.21823.

\bibitem{EwasiukProfumo2026}
C.~Ewasiuk and S.~Profumo,
\emph{Ghost degrees of freedom without quantum runaway: exact
moment bounds from an operator conservation law},
arXiv:2604.21348.

\bibitem{CJM2004}
J.~M.~Cline, S.~Jeon, and G.~D.~Moore,
\emph{The phantom menaced: constraints on low-energy effective
ghosts},
Phys.\ Rev.\ D \textbf{70}, 043543 (2004), hep-ph/0311312.

\bibitem{GV2013}
J.~Garriga and A.~Vilenkin,
\emph{Living with ghosts in Lorentz invariant theories},
JCAP \textbf{01} (2013) 036, arXiv:1202.1239.

\bibitem{Lagnese2021}
G.~Lagnese, F.~M.~Surace, M.~Kormos, and P.~Calabrese,
\emph{False vacuum decay in quantum spin chains},
Phys.\ Rev.\ B \textbf{104}, L201106 (2021), arXiv:2107.10176.

\bibitem{AbelSpannowskyWilliams2025}
S.~Abel, M.~Spannowsky, and S.~Williams,
\emph{Qumode tensor networks for false vacuum decay in quantum
field theory},
arXiv:2506.17388.

\bibitem{BurshteinGoldstein2024}
A.~Burshtein and M.~Goldstein,
\emph{Quantum simulation of the microscopic to macroscopic crossover
using superconducting quantum impurities},
Phys.\ Rev.\ B \textbf{111}, 174303 (2025), arXiv:2410.21941.

\end{thebibliography}

\end{document}
